\section{Identical particles}
\label{sec:q-identical}

Let the one-particle space have \(d\ge1\) specified orthonormal modes.
For \(N\) labelled particles a basis of the tensor power is indexed by
words \(w=(i_1,\ldots,i_N)\) in these modes; the empty word is the
vacuum. Permuting the slots by \(\pi\) defines a unitary \(U_\pi\),
with the convention \((\pi w)_s=w_{\pi^{-1}(s)}\).
For \(\chi=1\) or \(\chi=\operatorname{sgn}\), put
\[
               P_{N,\chi}=\frac1{N!}\sum_{\pi\in S_N}\chi(\pi)U_\pi .
\]
This rational matrix is a self-adjoint projection. Inversion preserves
\(\chi\), giving self-adjointness, and in the square every product
\(\tau\) occurs \(N!\) times with sign \(\chi(\tau)\).
Relabelling gives \(U_\tau P_{N,\chi}=\chi(\tau)P_{N,\chi}\).
Conversely, averaging fixes every vector with that transformation law.
The two choices specify the Bose and Fermi sectors, respectively; the
exchange law is a physical choice of sector. The finite construction may
test all of the \(N^N\) lists of length \(N>0\) for being a permutation
and count inversions for the sign. For \(N=0\), use one empty permutation.

\paragraph{Occupation coordinates.}
For bosons choose \(\alpha=(\alpha_1,\ldots,\alpha_d)\) in
\(\mathbb N^d\) with \(|\alpha|=\sum_j\alpha_j=N\).
There are \(m_\alpha=N!/\prod_j \alpha_j!\) words with these
multiplicities, by successive choices of slots. Their sum divided by
\(\sqrt{m_\alpha}\) is a unit vector \(|\alpha\rangle\).
These vectors are orthogonal because their supports are disjoint and
they span the symmetric sector because all coefficients on an orbit
are equal. For fermions a repeated letter forces the coefficient to
equal its negative under a transposition of equal slots. It therefore
vanishes. For each remaining increasing word take the signed sum of
its permutations divided by \(\sqrt{N!}\); on this orbit the
transformation law determines all coefficients from the increasing
word. This gives an orthonormal basis indexed by bit vectors
\(\alpha\in\{0,1\}^d\) of weight \(N\).
Consequently the sector dimensions are \({N+d-1\choose d-1}\) and
\({d\choose N}\), with the latter zero for \(N>d\). The first count
chooses the positions of \(d-1\) dividers in \(N+d-1\) slots; the
second chooses the occupied modes. Every denominator here has a positive
integer lower bound before taking its square root.

\paragraph{Creation, annihilation and domains.}
The Fock space is the Hilbert direct sum over \(N\) of these sectors.
Its effective vectors can be presented by finite occupation-coordinate
approximants with the usual \(\ell^2\) error witnesses. On finite
support set
\[
 \begin{aligned}
 a_j^*|\alpha\rangle&=\sqrt{\alpha_j+1}\,|\alpha+e_j\rangle,&
 a_j|\alpha\rangle&=\sqrt{\alpha_j}\,|\alpha-e_j\rangle
 &&\text{(bosons)} .
 \end{aligned}
\]
The term with \(\alpha_j=0\) in the annihilation formula is zero. For
fermions, creation is zero for an occupied mode and annihilation is zero
for an empty mode; otherwise change that bit and multiply by
\((-1)^{\sum_{k<j}\alpha_k}\). The two matrices are adjoints on the
finite-support domain by transposing each allowed basis transition.
For bosons these formulas give \([a_i,a_j^*]=\delta_{ij}I\) and
\([a_i,a_j]=[a_i^*,a_j^*]=0\). Indeed, on the same mode the two
number-preserving products have coefficients \(\alpha_i+1,\alpha_i\);
on different modes the two square-root factors are independent and the
same forbidden moves give zero in both orders. For fermions the
three formulas hold with anticommutators \(ST+TS\) in place of
commutators. On the same bit the two number-preserving products are its
empty and occupied projections and each doubled creation or annihilation
vanishes. On different bits the move at the smaller index flips the
sign for the move at the larger index while the reverse sign is unchanged.
An impossible move again gives zero in both orders.

For bosons the sector norms of creation and annihilation are bounded by
\(\sqrt{N+1}\) and \(\sqrt N\) because distinct admissible basis
vectors have distinct images and the squared coefficients satisfy these
bounds squared. On Fermi Fock space the corresponding bounds are one.
Thus an explicit sufficient domain witness for both Bose operations
is a cutoff \(M(p)\) with
\[
           \sum_{|\alpha|>M(p)} (|\alpha|+1)|c_\alpha|^2\le2^{-2p}.
\]
Each tail statement can equivalently be tested on every finite partial
sum. Apply the finite matrix to sectors through \(M(p+1)\) and allocate
the remaining half-error to the coefficients. For example take input
norm error at most \(2^{-p-1}/(M(p+1)+1)\) on this finite truncation.
The unboundedness is visible on states with \(N\) particles in mode
\(j\). A supplied one-particle unitary \(U\) acts on each sector by
\(U^{\otimes N}\), which commutes with permutations and is an isometry
with inverse \((U^*)^{\otimes N}\). It therefore acts on the direct sum
with the original tail bound.

For real mode energies \(\omega_j\), the free Hamiltonian on the
occupation basis multiplies by \(E_\alpha=\sum_j\omega_j\alpha_j\)
in units with \(\hbar=1\). It equals
\(\sum_j\omega_j a_j^*a_j\) on finite support, by the basis formulas.
Its unitary evolution multiplies each coefficient by
\(\exp(-itE_\alpha)\), preserving all unweighted tails and satisfying
the group law coordinatewise. For an effective derivative at time
\(t\), supply a cutoff for tails of
\(\sum_\alpha |E_\alpha c_\alpha|^2\). The tail error between the
difference quotient and \(-iE_\alpha\exp(-itE_\alpha)c_\alpha\) is at
most twice the norm of that weighted tail, since
\(|e^{ix}-1|\le |x|\). On a finite cutoff with \(|E_\alpha|\le K\),
the error is at most \(|h|K^2\|c\|_2/2\) for step \(h\), from
\(|e^{ix}-1-ix|\le x^2/2\) (integrate \(|e^{iy}-1|\le |y|\) on the
real segment). These are explicit input-dependent moduli for
Schr\"odinger's equation on the supplied domain.

\section{Open-system dynamics}
\label{sec:q-open}

For a finite-dimensional system an explicit time-independent
Markovian generator consists of a Hermitian matrix \(H\) and finitely
many matrices \(L_1,\ldots,L_r\). Write
\[
 \begin{aligned}
 M&=\sum_j L_j^*L_j, & \mathcal D(X)&=\sum_j L_jXL_j^*-\tfrac12(MX+XM),\\
 \mathcal H(X)&=-i(HX-XH), & \mathcal G&=\mathcal H+\mathcal D.
 \end{aligned}
\]
The following argument identifies the whole finite-dimensional class
when an effective right derivative at zero is part of the semigroup
data. For superoperators use the estimate
\(\|\mathcal A(X)\|\le b\|X\|\) with matrix operator norm estimates;
every use can equivalently be read as a bounded-vector test.

\begin{theorem}[Finite Lindblad evolution]\label{thm:q-lindblad}
For \(t\ge0\), \(S_t=\exp(t\mathcal G)\) is a channel,
\(S_{t+s}=S_t S_s\), and \(d S_t/dt=\mathcal G S_t\) (from the right
at zero). On a bounded interval it admits a primitive-recursive
approximation by channels that preserve trace at every step.
\end{theorem}
\begin{proof}
Choose integers \(a\ge2\|H\|\), \(\ell\ge\sum_j\|L_j\|^2\),
\(C=1+a+2\ell\), and \(T\ge\max(1,t)\).
For \(h=t/n\), \(n\ge CT\), take
\[
 K_0=(I-hM)^{1/2},\quad K_j=\sqrt h\,L_j,\quad
 U_h=\exp(-ihH),\qquad
 F_h(X)=U_h\left(\sum_{j=0}^{r}K_j X K_j^*\right)U_h^* .
\]
The input \(I-hM\) is positive since \(0\le M\le\ell I\). These
are Kraus matrices followed by a unitary, and
\(\sum_{j=0}^r K_j^*K_j=I\), so every \(F_h\) is a channel.
The binomial series of Theorem~\ref{thm:q-sqrt} gives
\[
 \|K_0-I\|\le h\ell,\qquad
 \|K_0-I+hM/2\|\le h^2\ell^2 .
\]
Indeed use \(F_m(hM)\) from that proof, whose nonconstant absolute
coefficients sum to at most one and whose linear coefficient is
\(-1/2\). Expanding \(K_0XK_0\) about \(X\) shows that the channel
inside parentheses differs from \(I+h\mathcal D\) by at most
\(3h^2\ell^2\) in superoperator norm. Here and below the identity
superoperator is denoted by \(I\) when it acts on matrices.

Unitary conjugation \(W_h(X)=U_hXU_h^*\) equals \(e^{h\mathcal H}(X)\):
expand the two exponentials and use the binomial formula for the
commuting operations of left and right multiplication. The same
identity gives unitarity by multiplication of the one-matrix series.
For \(0\le hC\le1\), the factorial series gives
\(\|W_h-I\|\le3ha\) and
\(\|W_h-I-h\mathcal H\|\le2h^2a^2\).
As \(\|\mathcal D\|\le2\ell\), multiplying the two approximations
bounds \(\|F_h-I-h\mathcal G\|\) by
\(3h^2\ell^2+2h^2a^2+6h^2a\ell\le 8h^2 C^2\).
The factorial tail of \(e^{h\mathcal G}\) beyond its linear part is
at most \(2h^2 C^2\). Also \(\|F_h\|\le1+h\ell\le e^{hC}\), since
\(K_0\) is a contraction, and \(\|e^{h\mathcal G}\|\le e^{hC}\).
Telescoping the difference of the \(n\)-th powers yields
\[
             \|F_{t/n}^{\,n}-e^{t\mathcal G}\|
               \le 10 C^2T^2\,4^{CT}/n .
\]
The factor \(4^{CT}\) bounds \(e^{CT}\) since \(C,T\) are integers.
Taking \(n\ge CT+10C^2T^2 4^{CT}2^{p+1}\) leaves half the precision
budget for finite matrix arithmetic. The limit has positive Choi
matrix by entry convergence on the finitely many matrix units, and
preserves trace by the same test. Thus it is a channel by
Theorem~\ref{thm:q-choi}.

For explicit series control at \(\|A\|\le D\), \(D\ge1\) integer,
let \(k=2D\). Beyond term \(k\) consecutive absolute terms have ratio
at most \(1/2\); term \(k\) is at most \(4^D\) by the scalar
exponential bound. Truncating after index \(k+p+2D+2\) therefore
makes the tail less than \(2^{-p}\). This estimate justifies
Cauchy products by finite truncation. The binomial identity for
commuting exponents gives the semigroup law; the linear Taylor
remainder bound above, with \(\mathcal G\) in place of \(\mathcal H\),
gives the derivative at each time after multiplication by \(S_t\).
\end{proof}

\begin{theorem}[Recovery of the generator]\label{thm:q-lindblad-converse}
Suppose \(T_0=I\) and \(T_t\) is a family of channels in dimension \(d\)
with a named right derivative \(G\) at zero and a primitive-recursive
accuracy modulus on each matrix unit. Then \(G\) has the above Lindblad
form with at most \(d^2\) operators \(L_j\). If the family is a semigroup,
\(T_t=e^{tG}\).
\end{theorem}
\begin{proof}
For a linear map \(A\) use its Choi matrix \(J_A\) from
Theorem~\ref{thm:q-choi}. Set \(w=\sum_i e_i\otimes e_i\),
\(P=ww^*/d\), \(Q=I-P\). Since \(J_I=ww^*\) and \(Qw=0\),
\[
       Q J_G Q=\lim_{h\downarrow0} Q J_{T_h}Q/h\ \ge0 .
\]
Take the limit along positive dyadics with the given modulus.
All \(J_{T_h}\) are Hermitian, so \(J_G\) is Hermitian; also
\(\operatorname{tr}G(X)=0\) by differentiating trace preservation.
Factor the positive matrix \(QJ_GQ\) with its positive square root
and reshape its \(d^2\) columns into matrices \(L_j\), as in the
Choi proof. Their map has Choi matrix \(QJ_GQ\).
Set
\[
             v=J_Gw/d-\frac{w^*J_Gw}{2d^2}w,
             \qquad V_{ai}=v_{ia}.
\]
Multiplication gives
\(J_G-QJ_GQ=J_GP+PJ_G-PJ_GP=vw^*+wv^*\). Testing matrix units,
the latter is the Choi matrix of \(X\mapsto VX+XV^*\).
Thus \(G(X)=\sum_j L_jXL_j^*+VX+XV^*\). Its vanishing trace on
matrix units gives \(V+V^*=-M\). With
\(H=(V^*-V)/(2i)\) this is the displayed generator.

For uniqueness from the semigroup law one can use finite telescoping.
In dimension \(d\), estimates on all units within \(\delta/d^2\)
give a superoperator bound \(\delta\), since
\(|X_{ij}|\le\|X\|\). Hence for each rational \(0<\delta\le1\) the
derivative modulus and the exponential Taylor bound give a positive
dyadic cutoff on \(h\) with
\(\|T_h-e^{hG}\|\le h\delta\) and
\(\|T_h\|\le 1+h(\|G\|+1)\) (use the derivative estimate directly for
the latter). These two closed inequalities also hold when \(h=0\)
by \(T_0=I\). For a general nonnegative named \(h\) below cutoff, a
strict violation of a tested inequality would exclude \(h>0\).
Since \(h\ge0\), this implies \(h=0\), contradicting the same violation.
This establishes the closed inequalities without deciding equality
to zero. For \(0\le t\le T\), take \(n\) large enough that \(h=t/n\)
lies under the cutoff. The difference of powers is bounded by
\(T e^{T(b+1)}\delta\) with an integer \(b\ge\|G\|\). Thus
\(T_t=e^{tG}\), with the required \(n\) obtained from the supplied
derivative modulus and the bound on \(T\). The modulus here applies to
every sufficiently small positive time, as in the definition of a
right derivative.
\end{proof}

\begin{example}[Spontaneous decay of one excited level]
In the ordered basis \(g,e\), take \(H=\omega |e\rangle\langle e|\)
and \(L=\sqrt\gamma\,|g\rangle\langle e|\) with \(\gamma\ge0\).
Set \(z=e^{-(\gamma/2+i\omega)t}\), \(q=|z|^2=e^{-\gamma t}\).
The evolved density matrix is
\[
 \begin{pmatrix}
 \rho_{gg}+(1-q)\rho_{ee} & \overline z\,\rho_{ge}\\
 z\rho_{eg} & q\rho_{ee}
 \end{pmatrix}.
\]
It follows either by multiplication with Kraus matrices
\(K_0=|g\rangle\langle g|+z|e\rangle\langle e|\) and
\(K_1=\sqrt{1-q}\,|g\rangle\langle e|\), or by differentiating:
the excited population solves \(x'=-\gamma x\), the ground population
gains \(\gamma x\), and the \(ge\) entry has coefficient
\(i\omega-\gamma/2\). The formula composes in time by multiplication
of \(z\) and \(q\), and agrees with the identity at zero. It includes
vanishing decay rate and vanishing initial excited population with the
same formulas.
\end{example}
